\documentclass[sigplan,10pt]{acmart}

\setcopyright{none}
\acmYear{2026}
\acmConference[ATC '26]{2026 ACM Annual Technical Conference}{2026}{}
\acmBooktitle{2026 ACM Annual Technical Conference}
\acmDOI{}
\acmISBN{}
\settopmatter{printfolios=true,printacmref=false}
\renewcommand\footnotetextcopyrightpermission[1]{}
\pagestyle{plain}

\usepackage{booktabs}
\usepackage{array}

\newcommand{\system}{GPUTriage}

\begin{document}
\sloppy

\title{LLM-Based Triage of GPU Numerical Failures from Public Issue Reports}

\author{Paper \#XXX}
\affiliation{\institution{Anonymous Submission}\country{Anonymous}}

\begin{abstract}
GPU software stacks increasingly combine compiler transformations, generated
kernels, mixed precision, and high-level Python interfaces. Their failure
reports are difficult to triage because symptoms such as NaNs, wrong answers,
overflow, dtype errors, compiler crashes, and performance regressions often
appear in the same public issue threads. This paper presents \system{}, a
human-labeled benchmark and LLM-based triage system for GPU numerical-failure
reports mined from public repositories. We separate weak retrieval labels from
human gold labels and build a 1,191-issue benchmark over eight GPU-related
software projects with a seven-class taxonomy. Weak keyword-derived labels
reach only 50.5\% accuracy, showing that issue-search labels are not reliable
ground truth. A fine-tuned FLAN-T5-base triage model reaches 80.5\% held-out
accuracy and 0.778 macro F1, while an LLM/deterministic agreement mode reaches
89.3\% accuracy at 68.3\% coverage. Leave-one-repository-out evaluation drops
to 66.8\%, exposing cross-project transfer as the central systems challenge.
We provide an anonymized artifact package with labels, audit logs, evaluation
tables, model outputs, and root-cause evidence files for review.
\end{abstract}

\maketitle

\section{Introduction}

GPU systems failures are no longer isolated to hand-written CUDA kernels. A
modern AI or scientific workload can pass through tensor libraries, graph
rewriters, JIT compilers, code generators, mixed-precision policies, and
vendor-specific execution backends before reaching hardware. When the result is
wrong, slow, non-finite, or impossible to compile, the first structured record
is often a public issue report. These reports are valuable systems data, but
they are noisy: a thread mentioning ``NaN'' may be a true invalid-value bug, a
missing-data question, a tolerance dispute, or a performance benchmark; a
``dtype'' issue may be a casting bug, unsupported API request, crash, or
precision error.

This paper asks whether public issue reports can support reliable triage of GPU
numerical failures. The task is practical for systems builders: maintainers
need to route reports to compiler, runtime, dtype, numerical-accuracy, or
performance experts; researchers need reusable traces of reliability problems;
and automated debugging tools need labeled reports that distinguish symptoms
from unrelated noise. Existing public issue trackers contain the raw material,
but weak labels from search queries are not enough.

We present \system{}, an LLM-based triage system and benchmark for categorizing
GPU numerical-failure reports. \system{} mines public GPU software issue
trackers, repairs weak labels through human annotation, and evaluates both
full-coverage and selective triage. The benchmark contains 1,191 human-labeled
issues from JAX, PyTorch, RAPIDS cuML, Triton, RAPIDS cuDF, CuPy, Numba, and
Apache TVM. Each issue receives one primary label from a seven-class taxonomy:
invalid values, overflow/underflow, precision/tolerance, dtype/casting,
crash/compile, performance-only, or not a numerical failure.

The contributions are:
\begin{itemize}
    \item We show that weak issue-search labels are insufficient for GPU
    numerical-failure triage, reaching only 50.5\% accuracy on the final
    benchmark.
    \item We construct a 1,191-issue human-labeled benchmark with audit logs,
    taxonomy repairs, a 250-row evidence-coded root-cause extension, and
    reproducibility artifacts.
    \item We implement and evaluate an LLM-based triage model that reaches
    80.5\% held-out accuracy and a selective agreement mode that reaches
    89.3\% accuracy at 68.3\% coverage.
    \item We evaluate chronological and leave-one-repository-out splits,
    identifying cross-project transfer as the main remaining barrier.
\end{itemize}

We organize the evaluation around four systems questions:
\textbf{RQ1}: Are weak labels from issue search usable as ground truth?
\textbf{RQ2}: Does a task-specific LLM improve report triage over transparent
text baselines? \textbf{RQ3}: Can the system trade coverage for higher
precision when reports are ambiguous? \textbf{RQ4}: Does triage transfer across
GPU software ecosystems, or does repository-specific language dominate?

\section{Problem and Dataset}

\subsection{Task}

Given a public issue report from a GPU software project, the triage task is to
predict the dominant failure category. The input includes title, body, repository
metadata, and available public issue context. The output is exactly one primary
label. This paper does not claim to prove root cause from text alone. Instead,
it classifies the report that a maintainer or triage assistant would see before
full debugging.

\subsection{Mining Public Reports}

We began with a silver seed mined from GPU-related repositories using terms such
as \texttt{nan}, \texttt{inf}, \texttt{overflow}, \texttt{dtype},
\texttt{precision}, \texttt{tolerance}, \texttt{incorrect}, and
\texttt{wrong result}. The retrieval query and keyword heuristics produced a
candidate label for each issue, but those labels were treated as weak
supervision rather than ground truth. We then constructed a pilot gold set and
expanded it with 1,000 additional human-labeled records. The final benchmark
contains 1,191 issues.

The final benchmark is deliberately repository-diverse rather than dominated by
one project. It includes compiler-centered reports, tensor-library reports,
Python/JIT reports, and GPU data-frame or array-library reports. That diversity
matters because the same surface symptom can mean different things across
systems. A dtype report in an array library often concerns API compatibility or
casting semantics; a dtype report in a compiler stack may indicate lowering,
legalization, or backend code generation. The benchmark therefore tests both
within-distribution triage and transfer across software ecosystems.

\begin{table}[t]
\centering
\caption{Human-labeled benchmark distribution.}
\label{tab:labels}
\begin{tabular}{lr}
\toprule
Primary label & Issues \\
\midrule
precision\_tolerance & 368 \\
dtype\_casting & 193 \\
crash\_compile & 192 \\
not\_numerical\_failure & 147 \\
nan\_inf & 143 \\
overflow\_underflow & 94 \\
performance\_only & 54 \\
\midrule
Total & 1191 \\
\bottomrule
\end{tabular}
\end{table}

\subsection{Taxonomy}

The taxonomy separates the observed report category from deeper root cause.
\texttt{nan\_inf} covers non-finite outputs; \texttt{overflow\_underflow}
covers numerical range failures; \texttt{precision\_tolerance} covers
wrong-answer and tolerance failures; \texttt{dtype\_casting} covers dtype
promotion, unsupported dtype, and casting semantics; \texttt{crash\_compile}
covers compiler/runtime exceptions and illegal memory failures;
\texttt{performance\_only} covers speed or memory complaints without a
correctness failure; and \texttt{not\_numerical\_failure} covers feature
requests, documentation issues, API questions, and unrelated reports.

\subsection{Human Label Quality}

The expanded annotation file included labels outside the final taxonomy,
including API feature requests, generic needs-review labels, and compiler/code
generation descriptions. We ran a deterministic taxonomy-normalization pass and
preserved the full repair log. In total, 180 rows were mapped into the final
taxonomy without silently dropping records. A blind audit over 119 rows produced
80.7\% agreement and Cohen's $\kappa=0.721$ between two annotators. A third
adjudication pass and a second 300-row adjudication pass focused on
low-confidence and model-error examples. Together, the passes applied 419 human
updates and changed 211 primary labels.

The review process also exposed weaknesses in the label schema. Early labels
mixed symptoms, suspected causes, and issue-management categories. We corrected
that by requiring a single primary report label and moving deeper cause evidence
into separate fields. For a maintainer-facing queue, this is the right
separation: first decide whether the report is numerical correctness, dtype,
compiler/runtime, performance, or non-numerical support; then ask whether later
evidence identifies a root cause.

\section{Triage System}

\system{} has four stages. First, the ingestion layer normalizes public issue
metadata and text into benchmark rows. Second, the taxonomy layer maps weak and
human labels into the final seven-class label space while preserving repair
provenance. Third, deterministic baselines train on TF-IDF, BM25, naive Bayes,
linear SVM, and logistic regression features. Fourth, the LLM stage fine-tunes a
FLAN-T5-base sequence-to-sequence classifier to emit one taxonomy label from
issue text.

\begin{table}[t]
\centering
\caption{Triage system stages.}
\label{tab:pipeline}
\begin{tabular}{p{0.27\columnwidth}p{0.60\columnwidth}}
\toprule
Stage & Role \\
\midrule
Mining & Retrieve public reports with numerical, dtype, compiler, and
wrong-answer terms. \\
Repair & Normalize weak and human labels into the seven-class taxonomy while
logging every repair. \\
Training & Train deterministic text baselines and a fine-tuned FLAN-T5-base
label generator. \\
Deployment & Emit a label, confidence mode, or abstention for human triage. \\
\bottomrule
\end{tabular}
\end{table}

The primary LLM split contains 945 training rows, 123 validation rows, and 123
held-out test rows. To reduce class imbalance, the training set is balanced by
oversampling minority classes to 2,058 examples. The selected checkpoint starts
from a prior expanded-gold checkpoint and is retrained on the v2 canonical
labels. We also evaluate selective modes. A deterministic voting ensemble can
abstain unless at least $k$ voters agree. The LLM/deterministic agreement mode
answers only when the fine-tuned LLM and deterministic ensemble predict the same
label.

This keeps the LLM role narrow. The model does not synthesize new reports,
rewrite evidence, or infer private information. It maps observed issue text into
an auditable taxonomy and can be compared directly with retrieval and linear
baselines. The deterministic ensemble remains useful because it provides a
transparent fallback and a disagreement signal for selective prediction.

\subsection{Implementation and Artifact}

The implementation is intentionally simple enough to be reproducible. The data
pipeline stores mined issue records, canonical labels, audit files, and
taxonomy-repair logs as CSV and JSONL files. Deterministic baselines use
standard sparse text features and are rerunnable from the released scripts. The
LLM training packet is emitted as sequence-to-sequence examples whose targets
are exactly the seven taxonomy labels. Evaluation scripts regenerate the
reported tables, per-class summaries, abstention metrics, chronological split,
and leave-one-repository-out split.

The artifact is part of the contribution rather than only a supplement. It
contains the benchmark, completed blind-audit files, adjudication outputs,
model predictions, error appendix, and root-cause evidence extension. During
review, public repository and archival links are removed from the paper to
preserve anonymity; the camera-ready version can restore those links after
acceptance.

\section{Evaluation}

\subsection{Protocol}

We report accuracy and macro F1. Macro F1 is needed because the benchmark is
imbalanced: precision/tolerance reports are the largest class, while
performance-only reports are the smallest. We evaluate four settings:
full-coverage shuffled evaluation, chronological evaluation on later issues,
leave-one-repository-out transfer, and selective prediction with abstention.

\subsection{Full-Coverage Results}

Table~\ref{tab:fullcoverage} shows that weak candidate labels reach only 50.5\%
accuracy, confirming that raw issue-search labels are not reliable supervision.
Linear text models substantially improve with human labels, and a voting
ensemble reaches 73.8\% accuracy and 0.712 macro F1. The fine-tuned LLM reaches
80.5\% held-out accuracy and 0.778 macro F1 on its test split, outperforming
deterministic baselines on the same held-out rows. A local zero-shot Llama 3.2
3B comparison reaches only 31.7\% accuracy and 0.322 macro F1 on those rows.

\begin{table}[t]
\centering
\caption{Full-coverage triage results.}
\label{tab:fullcoverage}
\begin{tabular}{lcc}
\toprule
Model & Accuracy & Macro F1 \\
\midrule
Majority baseline & 0.309 & 0.067 \\
Candidate weak label & 0.505 & 0.426 \\
BM25 nearest neighbor & 0.571 & 0.534 \\
TF-IDF logistic regression & 0.715 & 0.688 \\
TF-IDF linear SVM & 0.732 & 0.703 \\
Voting ensemble & 0.738 & 0.712 \\
FLAN-T5-base LLM & \textbf{0.805} & \textbf{0.778} \\
\bottomrule
\end{tabular}
\end{table}

The weak-label result answers RQ1: keyword retrieval is useful for candidate
collection but not for supervision. The 50.5\% weak-label accuracy is only
better than the majority baseline because issue-search terms are correlated
with failure language; it still mislabels many reports whose text contains
generic stack traces, dtype names, or numerical terms without a numerical
failure. The LLM result answers RQ2: the task-specific model improves over
transparent deterministic baselines, but the deterministic models remain useful
for debugging and confidence estimation.

\subsection{Generalization and Selective Triage}

Chronological evaluation trains on older issues and tests on newer issues. The
best model in this setting, bigram TF-IDF logistic regression, reaches 76.6\%
accuracy and 0.738 macro F1, while the voting ensemble reaches 74.5\% accuracy.
Leave-one-repository-out evaluation is harder: the best deterministic models
reach 66.8\% accuracy. CuPy and Numba are the weakest held-out repositories
because dtype/API compatibility, CUDA runtime failures, and compiler language
blur the boundary between numerical and non-numerical categories.

\begin{table}[t]
\centering
\caption{Generalization beyond shuffled evaluation.}
\label{tab:generalization}
\begin{tabular}{lcc}
\toprule
Setting / model & Accuracy & Macro F1 \\
\midrule
Chronological, linear SVM & 0.732 & 0.697 \\
Chronological, bigram logistic & \textbf{0.766} & \textbf{0.738} \\
Chronological, voting ensemble & 0.745 & 0.722 \\
Leave-one-repo, linear SVM & \textbf{0.668} & \textbf{0.648} \\
Leave-one-repo, voting ensemble & \textbf{0.668} & 0.647 \\
\bottomrule
\end{tabular}
\end{table}

\begin{table}[t]
\centering
\caption{Weakest leave-one-repository-out transfer cases.}
\label{tab:repo_weakness}
\begin{tabular}{lcc}
\toprule
Held-out repo & Issues & Ensemble acc. \\
\midrule
Numba & 65 & 0.462 \\
CuPy & 82 & 0.439 \\
RAPIDS cuDF & 132 & 0.545 \\
\bottomrule
\end{tabular}
\end{table}

This is the most important negative result. Shuffled and chronological splits
show that the benchmark is learnable, but holding out a whole repository
exposes vocabulary and workflow shift. Numba issues often mix Python exceptions,
LLVM compilation, CUDA behavior, and user-support language. CuPy issues often
mix dtype compatibility with numerical symptoms. These are exactly the cases
where a systems triage assistant should surface uncertainty rather than force a
confident label.

\begin{table}[t]
\centering
\caption{Deterministic ensemble ablation.}
\label{tab:ablation}
\begin{tabular}{lcc}
\toprule
Mode & Accuracy & Macro F1 \\
\midrule
Full deterministic ensemble & 0.768 & 0.753 \\
No weak candidate label & 0.761 & 0.748 \\
Linear-model vote only & 0.757 & 0.745 \\
Candidate label + SVM & 0.765 & 0.751 \\
\bottomrule
\end{tabular}
\end{table}

Table~\ref{tab:ablation} shows that the weak candidate label is helpful as a
feature-like signal but not as truth. Removing it lowers the deterministic
ensemble only slightly, while using it alone is far worse. This supports the
design choice to preserve weak labels for analysis but base the benchmark on
human annotation.

\begin{table}[t]
\centering
\caption{Selective prediction with abstention.}
\label{tab:abstain}
\small
\begin{tabular}{@{}lccc@{}}
\toprule
Rule & Cov. & Acc. & F1 \\
\midrule
Vote $\geq$ 3 & 0.926 & 0.769 & 0.744 \\
Vote $\geq$ 4 & 0.735 & 0.842 & 0.803 \\
Vote $\geq$ 5 & 0.301 & 0.939 & 0.733 \\
LLM/ensemble agree & 0.683 & 0.893 & 0.844 \\
\bottomrule
\end{tabular}
\end{table}

\subsection{Error Analysis}

The dominant errors are boundary cases that systems maintainers also find hard.
The largest confusion is \texttt{dtype\_casting} predicted as
\texttt{precision\_tolerance}; many dtype failures appear first as numerical
mismatches. \texttt{nan\_inf} reports are often predicted as
\texttt{precision\_tolerance} when the issue is written around failed accuracy
checks. \texttt{crash\_compile} reports are confused with numerical labels when
the exception occurs inside a failed computation. These patterns motivate
multi-label secondary-cause prediction and deeper root-cause linking to fixes.

The qualitative appendix contains concrete issue snippets for these boundary
cases. These examples are important for ATC because they show that the benchmark
is not an artificial label-matching exercise. Reports often contain enough
evidence to route the issue but not enough evidence to prove root cause, and the
hard cases are exactly where maintainers would want a tool to defer.

\section{Artifact and Threats}

The review artifact includes the benchmark, annotation guide, blind-audit files,
adjudication files, taxonomy repair logs, model predictions, evaluation tables,
error appendix, training packet, and root-cause evidence extension. Public
identity-bearing repository and archival links are intentionally omitted from
the blinded paper; an anonymized artifact package can accompany the ATC
submission.

The root-cause extension is not used to inflate the primary triage claim. It
contains a 50-row human-adjudicated root-cause subset and a 250-row
evidence-coded extension that records linked-fix signals, file categories, and
evidence provenance. We include it because the extension supports the next
systems question after triage: whether report symptoms can be connected to
fixes in compiler, runtime, dtype, test, documentation, or kernel code.

\begin{table}[t]
\centering
\caption{Root-cause extension label counts.}
\label{tab:rootcause}
\begin{tabular}{lr}
\toprule
Evidence label & Rows \\
\midrule
Numerical algorithm / tolerance & 74 \\
Non-bug, feature, docs, or support & 69 \\
Dtype or type semantics & 50 \\
Compiler backend or runtime & 30 \\
Performance or resource behavior & 27 \\
\midrule
Total & 250 \\
\bottomrule
\end{tabular}
\end{table}

Threats to validity remain. The benchmark is public-issue based and may not
represent private industrial failures. The full 1,191-row benchmark is not
double-adjudicated, although targeted audits and adjudication passes were used
to repair label drift. Some reports lack minimal reproducers or confirmed fixes,
so the task is report triage rather than definitive root-cause attribution.
Cross-repository transfer remains weaker than shuffled evaluation and should be
treated as an open systems problem.

\section{Related Work}

\system{} is related to three lines of work. First, GPU compiler and runtime
systems such as Triton, TVM, PyTorch, CuPy, and Numba expose the failure
surfaces studied in this paper. Second, software repository mining uses issue
reports, bug reports, and pull requests as empirical software-engineering data;
our contribution is a task-specific benchmark for GPU numerical-failure triage
rather than a general issue-classification dataset. Third, LLMs and
instruction-tuned sequence models are increasingly used for software tasks, but
our evaluation emphasizes human label quality, weak-label failure, and
cross-repository transfer rather than a single shuffled accuracy number.

\section{Conclusion}

\system{} turns noisy public GPU issue reports into a reusable benchmark and
LLM-based triage system for numerical failures. Human labels improve strongly
over weak search labels, a fine-tuned LLM reaches 80.5\% held-out accuracy, and
selective LLM/deterministic agreement reaches 89.3\% accuracy at 68.3\%
coverage. The remaining weakness is cross-project transfer, where repository
specific language and failure modes reduce accuracy to 66.8\%. The result is a
practical artifact for systems reliability research and a foundation for future
root-cause prediction over GPU software failures.

\begin{acks}
Omitted for double-blind review.
\end{acks}

\bibliographystyle{ACM-Reference-Format}
\begin{thebibliography}{9}

\bibitem{triton2019}
Tillet, P., Kung, H., and Cox, D. 2019. Triton: An intermediate language and
compiler for tiled neural network computations. In \emph{MAPL}.

\bibitem{cupy2017}
Okuta, R., Unno, Y., Nishino, D., Hido, S., and Loomis, C. 2017. CuPy: A
NumPy-compatible library for NVIDIA GPU calculations. In \emph{Workshop on
Machine Learning Systems at NeurIPS}.

\bibitem{pytorch2019}
Paszke, A. et al. 2019. PyTorch: An imperative style, high-performance deep
learning library. In \emph{NeurIPS}.

\bibitem{numba2015}
Lam, S. K., Pitrou, A., and Seibert, S. 2015. Numba: A LLVM-based Python JIT
compiler. In \emph{LLVM-HPC}.

\bibitem{tvm2018}
Chen, T. et al. 2018. TVM: An automated end-to-end optimizing compiler for deep
learning. In \emph{OSDI}.

\bibitem{t5}
Raffel, C. et al. 2020. Exploring the limits of transfer learning with a
unified text-to-text transformer. \emph{JMLR} 21, 140, 1--67.

\bibitem{flan}
Chung, H. W. et al. 2024. Scaling instruction-finetuned language models.
\emph{JMLR} 25, 70, 1--53.

\end{thebibliography}

\end{document}
